% Part: model-theory
% Chapter: interpolation
% Section: definability

\documentclass[../../../include/open-logic-section]{subfiles}

\begin{document}

\olfileid{mod}{int}{def}

\section{વ્યાખ્યેયતા પ્રમેય}

અંતર્વેશન પ્રમેયનો એક મહત્વનો ઉપયોગ બેથનું વ્યાખ્યેયતા પ્રમેય છે.
$n$-સ્થાની સંબંધ~$R$ને વ્યાખ્યાયિત કરવા માટે આપણે $R$ ન આવતો હોય
અને $n$ મુક્ત !!{variable}s ધરાવતો કોઈ !!a{formula}~$!C$ આપી શકીએ
છીએ. આ $!C$ દ્વારા $R$ની \emph{સ્પષ્ટ} વ્યાખ્યા થાય. પછી આપણે એમ
પણ કહી શકીએ કે $n$-સ્થાની !!{predicate}~$P$ ધરાવતી કોઈ !!{language}માં
સિદ્ધાંત~$\Sigma(P)$, $P$ને સ્પષ્ટ રીતે વ્યાખ્યાયિત કરે છે જો તેમાં
ઔપચારિક કરેલી સ્પષ્ટ વ્યાખ્યા હોય (અથવા ઓછામાં ઓછું તેમાંથી ફલિત થતી
હોય), એટલે કે,
\[
\Sigma(P) \Entails \lforall[x_1][\dots
  \lforall[x_n][(\Atom{P}{x_1,\dots, x_n} \liff !C(x_1, \dots,
    x_n))]].
\]
પરંતુ સંબંધને વ્યાખ્યાયિત કરવાનો---તેને સંપૂર્ણપણે નિશ્ચિત કરવાના
અર્થમાં---સ્પષ્ટ વ્યાખ્યા માત્ર એક માર્ગ છે. સિદ્ધાંત એવો પણ હોઈ શકે
કે કોઈ પણ નિદર્શમાં !!{language}ના બાકીના ભાગના અર્થઘટનથી $P$નું
અર્થઘટન નિશ્ચિત થાય. વ્યાખ્યેયતા પ્રમેય કહે છે કે જ્યારે પણ સિદ્ધાંત
$P$નું અર્થઘટન આ રીતે નિશ્ચિત કરે---જ્યારે તે $P$ને \emph{ગૂઢ રીતે
વ્યાખ્યાયિત કરે}---ત્યારે તે તેને સ્પષ્ટ રીતે પણ વ્યાખ્યાયિત કરે છે.

\begin{defn}
ધારો કે $\Lang{L}$ એવી !!a{language} છે જેમાં !!{predicate}~$P$ નથી.
$\Lang{L}\cup\{P\}$નાં !!{sentence}sનો ગણ~$\Sigma(P)$, $P$ને ત્યારે
અને માત્ર ત્યારે \emph{સ્પષ્ટ રીતે વ્યાખ્યાયિત કરે છે} જ્યારે
$\Lang{L}$નું એવું !!a{formula}~$!C(x_1,\dots,x_n)$ છે કે
\[
\Sigma(P) \Entails \lforall[x_1][\dots
  \lforall[x_n][(\Atom{P}{x_1,\dots, x_n} \liff !C(x_1, \dots,
    x_n))]].
\]
\end{defn}

\begin{defn}
ધારો કે $\Lang{L}$ એવી !!a{language} છે જેમાં !!{predicate}s $P$
અને~$P'$ નથી. $\Lang{L}\cup\{P\}$નાં !!{sentence}sનો ગણ~$\Sigma(P)$,
$P$ને ત્યારે અને માત્ર ત્યારે \emph{ગૂઢ રીતે વ્યાખ્યાયિત કરે છે}
જ્યારે
\[
\Sigma(P) \cup \Sigma(P') \Entails \lforall[x_1][\dots
  \lforall[x_n][(\Atom{P}{x_1,\dots, x_n} \liff \Atom{P'}{x_1,\dots,
      x_n})]],
\]
જ્યાં $\Sigma(P')$ એ $\Sigma(P)$માં $P$ને સર્વત્ર $P'$થી બદલીને
મળતું પરિણામ છે.
\end{defn}

બીજા શબ્દોમાં, કોઈ પણ નિદર્શ~$\Struct{M}$ અને $R,R'\subseteq
\Domain{M}^n$ માટે, જો $\Expan{M}{R}\Entails\Sigma(P)$ અને
$\Expan{M}{R'}\Entails\Sigma(P')$ બંને, તો $R=R'$; અહીં
$\Expan{M}{R}$ એ $\Lang{L}$ના $\Lang{L}\cup\{P\}$ સુધીના વિસ્તાર
માટેની એવી !!{structure}~$\Struct{M'}$ છે કે $\Assign{P}{M'}=R$, અને
$\Expan{M}{R'}$ માટે પણ એ જ રીતે.

\begin{thm}[બેથનું વ્યાખ્યેયતા પ્રમેય]
$\Lang{L}\cup\{P\}$નાં !!{sentence}sનો ગણ~$\Sigma(P)$, $P$ને ગૂઢ
રીતે વ્યાખ્યાયિત કરે છે જો અને માત્ર જો $\Sigma(P)$, $P$ને સ્પષ્ટ
રીતે વ્યાખ્યાયિત કરે છે.
\end{thm}
\sourcecorrection{OLINT-008}{સ્થિર અંગ્રેજી મૂળની
  \texttt{definability.tex}, પંક્તિઓ~66--68માં પ્રમેય $\Sigma(P)$ને
  સૂત્રોનો ગણ કહે છે, જ્યારે તરત અગાઉની બંને વ્યાખ્યાઓ તેને
  વાક્યોનો ગણ કહે છે અને નીચેની સાબિતી પણ સાન્ત ઉપગણનાં
  વાક્યોના સંયોજન પર ચાલે છે. અહીં વ્યાખ્યાઓ અને સાબિતી મુજબ
  વાક્યો પુનઃસ્થાપિત કર્યાં છે.}
\sourcecorrection{OLINT-009}{એ જ મૂળની પંક્તિ~67માં દ્વિમુખી શરતના
  ``if and only if''માંથી છેલ્લું ``if'' છૂટી ગયું છે. અહીં સંપૂર્ણ
  દ્વિમુખી શરત પુનઃસ્થાપિત કરી છે.}

\begin{proof}
જો $\Sigma(P)$, $P$ને સ્પષ્ટ રીતે વ્યાખ્યાયિત કરે, તો આ બંને સાચાં
છે:
\begin{align*}
  \Sigma(P) & \Entails & \lforall[x_1][\dots \lforall[x_n]
    [(\Atom{P}{x_1,\dots, x_n} \liff !C(x_1,\dots,x_n))]]\\
  \Sigma(P') & \Entails & \lforall[x_1][\dots \lforall[x_n]
    [(\Atom{P'}{x_1,\dots, x_n} \liff !C(x_1,\dots,x_n))]]
\end{align*}
અને નિષ્કર્ષ ફલિત થાય છે. વ્યસ્ત દિશા માટે ધારો કે $\Sigma(P)$,
$P$ને ગૂઢ રીતે વ્યાખ્યાયિત કરે છે. પહેલા $\Lang{L}$માં !!{constant}s
$c_1$,~\dots,~$c_n$ ઉમેરીએ. ત્યારે
\[
\Sigma(P) \cup \Sigma(P') \Entails
\Atom{P}{c_1, \dots, c_n} \to \Atom{P'}{c_1, \dots, c_n}.
\]
સઘનતાથી એવા સાન્ત ગણો $\Delta_0\subseteq\Sigma(P)$ અને
$\Delta_1\subseteq\Sigma(P')$ છે કે
\[
\Delta_0 \cup \Delta_1 \Entails
\Atom{P}{c_1, \dots, c_n} \to \Atom{P'}{c_1, \dots, c_n}.
\]
એવા બધા !!{sentence}s~$!A(P)$નું સંયોજન~$!D(P)$ કહો કે કાં તો
$!A(P)\in\Delta_0$ અથવા $!A(P')\in\Delta_1$, અને એવા બધા
!!{sentence}s~$!A(P')$નું સંયોજન~$!D(P')$ કહો કે કાં તો
$!A(P)\in\Delta_0$ અથવા $!A(P')\in\Delta_1$. ત્યારે
$!D(P)\land !D(P')\Entails\Atom{P}{c_1,\dots,c_n}\to
\Atom{P'}{c_1,\dots,c_n}$. આને એવી રીતે ફરી ગોઠવી શકાય કે દરેક
!!{predicate}, $\Entails$ની એક બાજુ જ આવે:
\sourcecorrection{OLINT-010}{સ્થિર અંગ્રેજી મૂળની
  \texttt{definability.tex}, પંક્તિ~96માં જમણી બાજુનું પરમાણુક સૂત્ર
  આસપાસનાં બધાં દાખલાથી જુદું $P'c_1\dots c_n$ લખાયું છે. અહીં એ જ
  $n$-સ્થાની પરમાણુક સૂત્ર માટેનું આવૃત્તિ-સુસંગત
  \texttt{\string\Atom\{P'\}\{c\_1,\string\dots,c\_n\}} સ્વરૂપ વાપર્યું
  છે.}
\[
!D(P) \land \Atom{P}{c_1, \dots, c_n} \Entails
!D(P') \to \Atom{P'}{c_1, \dots, c_n}.
\]
ક્રેગના અંતર્વેશન પ્રમેય મુજબ $P$ કે~$P'$ ન ધરાવતું એવું
!!a{sentence}~$!C(c_1,\dots,c_n)$ છે કે
\begin{align*}
  !D(P) \land \Atom{P}{c_1, \dots, c_n} & \Entails !C(c_1,\dots, c_n); \\
  !C(c_1,\dots, c_n) & \Entails !D(P') \to \Atom{P'}{c_1, \dots, c_n}.
\end{align*}
આ બે ફલિતતાઓમાંથી પહેલામાંથી મળે છે: $!D(P)\Entails
\Atom{P}{c_1,\dots,c_n}\lif !C(c_1,\dots,c_n)$. અને બીજામાંથી,
કારણ કે કોઈ $\Lang{L}\cup\{P\}$-નિદર્શ $\Expan{M}{R}\Entails
!A(P)$ જો અને માત્ર જો તેને અનુરૂપ $\Lang{L}\cup\{P'\}$-નિદર્શ
$\Expan{M}{R}\models !A(P')$, આપણને
$!C(c_1,\dots,c_n)\Entails !D(P)\lif\Atom{P}{c_1,\dots,c_n}$ મળે છે,
અને તેમાંથી
\[
!D(P) \Entails !C(c_1,\dots,c_n) \to \Atom{P}{c_1,\dots, c_n}.
\]
બંનેને સાથે મૂકતાં $!D(P)\Entails\Atom{P}{c_1,\dots,c_n}\liff
!C(c_1,\dots,c_n)$ મળે છે, અને એકદિશવર્ધિતા તથા સર્વત્રીકરણથી આ પણ
મળે છે:
\[
\Sigma(P) \Entails
\lforall[x_1][\dots\lforall[x_n][(\Atom{P}{x_1,\dots, x_n} \liff
    !C(x_1,\dots, x_n))]].
\]
\end{proof}

\end{document}
